\section{Introduction}
\label{sec:intro}
Multi-task benchmarks are the main infrastructure for comparing machine-learning models: a leaderboard is a table of scores on several sub-datasets, collapsed into one ranking by an aggregation rule, usually an average. Rankings steer adoption, publication and funding, so developers have a reason to act strategically. One cheap action is to include data related to some sub-datasets in the training of a chosen model, which raises its scores there~\citep{zhou2023cheater,yang2023rephrased}. How much of this does it take to reach first place, and does the answer depend on the aggregation rule?

We study this on public leaderboards. We formalise a \emph{targeted investment} (a target model raises its score on chosen sub-datasets up to a cap, all other models fixed), measure it by the number of sub-datasets touched ($K$), the total score points ($L_1$) and the largest single jump ($L_\infty$), and compute, \emph{exactly}, the cheapest investment that makes a given model strictly first under each of a menu of rules: mean-type rules, order-statistic rules (trimmed mean, median, minimum) and ordinal rules (Borda, plurality, Condorcet winner). Because the cap is an assumption rather than a measurement, we repeat everything under several cap families, one of them calibrated on observed fine-tune gains.

\paragraph{Contributions.}
(1) A formal model of targeted investment with three cost measures and an explicit cap assumption, and exact minimum-investment solvers for ten aggregation rules, verified against brute force; closed forms for the mean, minimum and (odd-$n$) median. (2) An empirical study on two public leaderboards (Open LLM Leaderboard, $M=4496$ models, $n=6$; HELM MMLU v1.13.0, $M=79$, $n=57$) with intervals over targets, sensitivity to caps, tie margin, evaluation noise, task resampling and data slices. (3) Findings (Section~\ref{sec:results}) that are not forced by the definitions: the order of rules by investment depends on the cost measure, the dataset and the cap, and which rule is easiest to topple under a binding cap flips with the cap; points are cheapest for the median, minimum and plurality, and expensive for averaging rules; targeting saves points mainly for non-mean rules, and for the mean only under binding caps. (4) Code, checksums and a script-to-result map (Appendix~\ref{app:repro}).

\paragraph{What this paper does not do.}
We do not train models, model rivals' reactions, or claim that any real developer behaved in this way; we also do not solve every rule on every target (ordinal rules use a sample) or compute Copeland, Schulze or Kemeny rankings.
